\documentclass[12pt]{article}
\usepackage{amsmath, amssymb}
\usepackage[margin=1in]{geometry}
\setlength{\parskip}{6pt}
\title{QUBO Formulation: Minimum Dominating Set Problem}
\date{}
\begin{document}
\maketitle

\subsection*{Minimum Dominating Set Problem}

\paragraph{Problem Statement.}
Given an undirected graph $G=(V,E)$ with vertex set $V$ and edge set $E$, the Minimum Dominating Set Problem seeks a subset $D \subseteq V$ of minimum cardinality such that every vertex $v \in V$ is either contained in $D$ or adjacent to at least one vertex in $D$. Formally, we wish to minimize $|D|$ subject to the domination condition.

\paragraph{Decision Variables.}
For each vertex $i \in V$, we introduce a binary decision variable $x_i \in \{0,1\}$. The variable $x_i$ takes the value $1$ if vertex $i$ is selected in the dominating set, and $0$ otherwise. The solution vector is $x = (x_i)_{i \in V} \in \{0,1\}^{|V|}$.

\paragraph{QUBO Derivation.}
We construct a Quadratic Unconstrained Binary Optimization (QUBO) Hamiltonian $H(x)$ by combining the objective function with penalty terms that enforce the domination constraints.

\textbf{Step 1.} The original objective is to minimize the number of selected vertices:
\begin{align}
    \min_{x \in \{0,1\}^{|V|}} \sum_{i \in V} x_i.
\end{align}
In the QUBO framework, this objective is directly encoded as the energy term $H_{\text{obj}}(x) = \sum_{i \in V} x_i$.

\textbf{Step 2.} The domination constraint requires that for every vertex $j \in V$, at least one vertex in its closed neighborhood $S_j = \{j\} \cup \{i \in V : (i,j) \in E\}$ is selected. This is expressed algebraically as:
\begin{align}
    \sum_{i \in S_j} x_i \ge 1, \quad \forall j \in V.
\end{align}

\textbf{Step 3.} To convert the inequality constraints into a penalty term, we penalize violations of the form $\sum_{i \in S_j} x_i < 1$. Since variables are binary, a violation occurs if and only if $\sum_{i \in S_j} x_i = 0$. We enforce the constraint using a quadratic penalty with weight $\lambda > 0$:
\begin{align}
    P_j(x) &= \lambda \left(1 - \sum_{i \in S_j} x_i\right)^2.
\end{align}
Expanding the square and applying the idempotent property $x_i^2 = x_i$ for binary variables yields:
\begin{align}
    P_j(x) &= \lambda \left(1 - 2\sum_{i \in S_j} x_i + \sum_{i \in S_j} x_i^2 + \sum_{i < k \in S_j} 2 x_i x_k\right) \\
    &= \lambda \left(1 - \sum_{i \in S_j} x_i + \sum_{i < k \in S_j} 2 x_i x_k\right).
\end{align}
Summing over all vertices gives the total penalty Hamiltonian:
\begin{align}
    H_{\text{pen}}(x) = \sum_{j \in V} P_j(x) = \sum_{j \in V} \lambda \left(1 - \sum_{i \in S_j} x_i + \sum_{i < k \in S_j} 2 x_i x_k\right).
\end{align}

\textbf{Step 4.} The complete QUBO Hamiltonian is the sum of the objective and penalty terms:
\begin{align}
    H(x) = H_{\text{obj}}(x) + H_{\text{pen}}(x) = \sum_{i \in V} x_i + \sum_{j \in V} \lambda \left(1 - \sum_{i \in S_j} x_i + \sum_{i < k \in S_j} 2 x_i x_k\right).
\end{align}

\textbf{Step 5.} Collecting like terms, we read off the QUBO matrix entries $Q_{i,k}$ and the constant offset $c$ in the standard form $H(x) = \sum_{i \le k} Q_{i,k} x_i x_k + c$. The general closed-form expressions are:
\begin{align}
    c &= \sum_{j \in V} \lambda = |V|\lambda, \\
    Q_{i,i} &= 1 - \lambda \cdot |\{j \in V : i \in S_j\}|, \quad \forall i \in V, \\
    Q_{i,k} &= 2\lambda \cdot |\{j \in V : i,k \in S_j\}|, \quad \forall i < k.
\end{align}
Here, $|\{j \in V : i \in S_j\}|$ equals the degree of vertex $i$, and $|\{j \in V : i,k \in S_j\}|$ counts the number of constraints in which both $x_i$ and $x_k$ appear.

\paragraph{Penalty Weight Justification.}
The penalty weight $\lambda$ must be chosen sufficiently large to ensure that any infeasible solution has strictly higher energy than the optimal feasible solution. Violating a single domination constraint corresponds to leaving a closed neighborhood completely unselected, which can increase the objective function by at most $|S_j| \le |V|$. Therefore, any $\lambda > |V|$ guarantees that the penalty for a single violation outweighs the maximum possible objective gain from removing a vertex. In practice, a conservative value such as $\lambda = 10$ is selected to provide a robust margin against numerical noise and to ensure that multiple simultaneous violations are heavily penalized.

\paragraph{Correctness Argument.}
Minimizing $H(x)$ over $\{0,1\}^{|V|}$ recovers the optimal dominating set due to the structure of the penalty formulation. First, any feasible solution satisfies all domination constraints, resulting in $H_{\text{pen}}(x) = 0$; thus, its energy equals the objective value $\sum_{i \in V} x_i$. Second, any infeasible solution violates at least one constraint, incurring a penalty of at least $\lambda$. Since $\lambda$ is chosen strictly greater than the maximum possible reduction in the objective function from any single variable flip, the energy of any infeasible configuration strictly exceeds that of the best feasible configuration. Consequently, the global minimum of $H(x)$ must be feasible and correspond to a minimum dominating set.

\paragraph{Illustrative Example.}
Consider a concrete instance with $n=3$ vertices, $V=\{0,1,2\}$, and a specific edge structure that yields the following closed-neighborhood overlaps. Using $\lambda=10$, we instantiate the general formulas from Step 5. The constant offset is $c = 3 \times 10 = 30$. The concrete objective is $\sum_{i=0}^2 x_i$. The concrete penalty terms expand to linear coefficients $Q_{i,i} = 1 - 10 \cdot \deg(i)$ and quadratic coefficients $Q_{i,k} = 20 \cdot |\{j \in V : i,k \in S_j\}|$. Substituting the instance-specific degrees and neighborhood overlaps yields $Q_{0,0}=-19$, $Q_{1,1}=-29$, $Q_{2,2}=-19$, $Q_{0,1}=40$, $Q_{0,2}=20$, and $Q_{1,2}=40$. The resulting Hamiltonian is:
\begin{align}
    H(x) &= 30 - 19x_0 - 29x_1 - 19x_2 + 40x_0x_1 + 20x_0x_2 + 40x_1x_2.
\end{align}
Evaluating $H(x)$ over all $2^3=8$ binary assignments reveals that the minimum energy is $10$, achieved at $x=(1,0,1)$ and $x=(0,1,0)$. The global minimum $x^*=(0,1,0)$ with energy $10$ correctly identifies vertex $1$ as the minimum dominating set for this instance.

\end{document}